\begin{defn}
A sequence $\{ x_n \}_{n=1}^\infty$ is said to \emph{converge} to a number
$x \in \R$ if for every $\epsilon > 0$, there exists an $M \in \N$ such
that $\sabs{x_n - x} < \epsilon$ for all $n \geq M$.
The number $x$ is called a \emph{limit} of the sequence.
We will prove that the limit $x$ is unique, and so we write%
\footnote{In text, this may get rendered as $\lim_{n\to\infty} x_n$.}
\glsadd{not:limseq}%
\begin{equation*}
\lim_{n\to \infty} x_n \coloneqq x .
\end{equation*}

A sequence
that converges is said to be \emph{convergent}\index{convergent!sequence}.
Otherwise, we say the sequence \emph{diverges}
or that it is
\emph{divergent}\index{divergent!sequence}.
\end{defn}
